\documentclass[12pt]{article}
\usepackage[utf8]{inputenc}
\usepackage[T1]{fontenc}
\usepackage{hyperref}
\title{Karami Language Resources for K-12 Curriculum Development}
\author{Compiled by Claude Sonnet 5.5 (claude-sonnet-5-5) for Andr\'as Kornai}
\date{}
\begin{document}
\maketitle

\section{Introduction}
Karami (ISO 639-3 xar; Glottocode kara1497) is an extinct language of southern Papua New Guinea (Western Province, near Daru). The sources found disagree on its classification. Ethnologue, as relayed by ASJP, places it in Trans-New Guinea (Inland Gulf, Minanibai). Pawley and Hammarst\"om (2018) treat it as an isolate. It is known only from a short word list, and the ASJP database gives 1959 as its extinction date.

\section{Primary Educational Resources}
No educational materials of any kind were found. None are listed here.

\section{Linguistic and Reference Materials}
The main source is a 1917 vocabulary by L.\ A.\ Flint, which places speakers in the villages of Kikimairi and Aduahai. Many loanwords come from Foia Foia.
\begin{itemize}
\item Wikipedia, ``Karami language''. \url{https://en.wikipedia.org/wiki/Karami_language}
\item ASJP database entry for Karami (word list and classification). \url{https://asjp.clld.org/languages/KARAMI}
\item Glottolog reference record for L.\ A.\ Flint, 1917--1918. \url{https://glottolog.org/resource/reference/id/109402}
\item Glottolog reference record for Karl J.\ Franklin, 1973 (classification). \url{https://glottolog.org/resource/reference/id/303435}
\item Glottolog reference record for Andrew Pawley, 1966. \url{https://glottolog.org/resource/reference/id/150153}
\item CrossGram entry for Karami (kara1497). \url{https://crossgram.clld.org/languages/kara1497}
\end{itemize}

\section{Religious and Community Resources}
No scripture or community organization related to Karami was found.

\section{Digital and Audio Resources}
No audio or digital teaching resources were found. The Endangered Languages Project has an entry for the language, which I could not check in detail. \url{https://endangeredlanguages.com/lang/10740?hl=de}

\section{Curriculum Development Context}
With only a short word list, Karami cannot support a language curriculum. At most, a K-12 unit could present the language as part of the history of the Western Province and of Papua New Guinea's linguistic diversity. Any such work should involve the local communities of the Daru area.
\end{document}